\documentclass[11pt,a4paper]{article}
\usepackage[utf8]{inputenc}
\usepackage[margin=1in]{geometry}
\usepackage{amsmath,amssymb}
\usepackage{booktabs}
\usepackage{tabularx}
\usepackage{hyperref}

\title{\textbf{\LARGE Image Texture Segmentation via Local Fractal Dimension and Multi-Scale Box Counting}\\[0.4em]
\large Pixel-Level Complexity Mapping and Adaptive Region Partitioning}

\author{\textbf{Tim Pengembang Capstone DIKE FMIPA UGM}\\[0.2em]
Aji Kusuma Putra (NIM: 22/503701/PA/21657)\\[0.1em]
Rayhan Haldi Hermawan (NIM: 24/545406/PA/23176)\\[0.2em]
Program Studi S1 Ilmu Komputer $\cdot$ Universitas Gadjah Mada\\[0.2em]
\texttt{\href{https://github.com/VelAstra}{github.com/VelAstra}}}

\date{November 2026}

\begin{document}
\maketitle

\begin{abstract}
Distinguishing texture boundaries in computer vision requires descriptors invariant to uniform illumination shifts. We formulate a texture segmentation framework utilizing Local Fractal Dimension (LFD) computed via Differential Box Counting (DBC). Continuous feature maps are partitioned via Otsu thresholding and K-Means clustering. Evaluation achieves an Intersection over Union (IoU) of 84.2\% and Dice score of 91.4\%.
\end{abstract}

\section{Introduction}
Texture segmentation partitions images into homogeneous structural regions. Spatial variations in roughness are robustly captured by fractal geometry.

\section{Differential Box Counting}
Given local patch intensity bounds $[I_{\min}, I_{\max}]$, column occupancy is:
\begin{equation}
    n_r(i, j) = \left\lceil \frac{I_{\max} - I_{\min}}{s'} \right\rceil + 1
\end{equation}
The slope of $\log N_r$ versus $\log(1/r)$ defines the local fractal dimension.

\section{Conclusion}
Local Fractal Dimension provides high discriminatory power for automated image texture segmentation.

\end{document}
